\clearpage
\fancyhead[L]{\small MOLLWEIDE. COMMENTATIONES.}
\hypertarget{native-572-lower}{}
\begin{center}\small Göttingische gelehrte Anzeigen. No. 71. Pages 705..710. 2 May 1814.\\[3pt]\rule{.72\linewidth}{.3pt}\end{center}

\emph{Commentationes mathematico-philologicae tres, sistentes explicationem duorum locorum difficilium, alterius Virgilii, alterius Platonis itemque examinationem duorum mensurarum praeceptorum Columellae. Adjecta est epistola ad v. cl.} J. G. \textsc{Schneider} \emph{de excerptis geometricis Epaphroditi et Vitruvii Rufi scripta ab auctore harum commentationum} \textsc{Carolo Brandano Mollweide}, astron. in acad. Lipsiensi professore. 1813. Leipzig. Published by Cnobloch. 122 pages in octavo, with one copperplate.

The rarer it is to find thorough mathematical knowledge and a deep familiarity with the writers and languages of antiquity in one person, the more welcome must be the light that the learned author of these essays sheds, with the help of mathematics, on several obscure passages in ancient writers. The first essay concerns the well-known passage in Virgil’s poem on agriculture about the times of the two honey harvests.

\SourceLeaf{573}\begin{quote}\small Bis gravidos cogunt fetus: duo tempora messis \\ Taygete simul os terris ostendit honestum \\ Pleias, et oceani spretos pede repulit amnes, \\ Aut eadem sidus fugiens ubi piscis aquosi \\ Tristior hibernas coelo descendit in undas.

\hfill Georg. IV. 231---235. \end{quote}

Most commentators agree that the season indicated in the last two lines is autumn, when the Pleiades set in the morning, and that descendit in undas must therefore refer to the \emph{cosmical} setting. But many different explanations of the words sidus fugiens piscis aquosi have been tried without any gaining general approval, and some commentators have admitted outright that the words as they stand are inexplicable. Mr.\ \textsc{Mollweide} set out his view of this passage some time ago in the Monatl.\ Corresp. (May 1802). Later, on his appointment as professor of astronomy in Leipzig, he published a revised version of that essay in a short work issued separately. It now appears in this collection with further additions. By piscis aquosus most commentators mean the fishes of the zodiac, whereas he means the Southern Fish. In his later revisions, however, he dropped one of the main arguments he had used to support this view in the first essay (namely, that the epithet aquosus, which seemed somewhat superfluous for the fishes of the zodiac, referred to the position of the first-magnitude star in the Southern Fish at the stream of water from Aquarius). Presumably he did so because there are other passages where piscis with the same epithet clearly means the fish of the zodiac. He now takes the epithet only as an allusion to the wet season. He explains the \emph{flight} of the Pleiades from the Southern Fish by the fairly close coincidence of the cosmical setting of the Pleiades and the acronychal rising of the Southern Fish. To prove that an image of flight can be used in this sense, he refers to another passage in the same poem.

\begin{quote}\small Vere fabis satio; tum te quoque, Medica, putres \\ Accipiunt sulci, et milio venit annua cura, \\ Candidus auratis aperit quum cornibus annum \\ Taurus, et adverso \emph{cedens} canis occidit astro.

\hfill Georg. I. 215---218. \end{quote}

\noindent Here cedens has the same meaning as fugiens in the passage above and, according to Mr.\ \textsc{Mollweide}, points to the close coincidence of the Dog’s evening setting and the Bull’s morning rising. For all the erudition with which Mr.\ \textsc{Mollweide} supports his explanation, the reviewer cannot deny that it still feels somewhat forced to him. He is quite willing to accept this explanation of the second passage. When the evening setting of one star coincides with the morning rising of another, and this is called the flight of the first from the second, he sees in it a beautiful and true poetic image \emph{despite the different} times of day. Yet it seems awkward to him to apply the same image to the reverse case (the one at issue here), where the fleeing star sets in the morning and the one driving it away rises in the evening. In spring the constellation of the Great Dog, which until then had shone in the western sky every evening, begins to disappear from view altogether at about the time when the Bull, which in turn had been quite invisible for a while, begins to reappear and increasingly adorn the eastern sky in the early morning hours. What objection can there be if a \SourceLeaf{574}poet calls this the flight of the Dog from the Bull? But what becomes of the truth and vividness of the poetic image when the poet has the Pleiades flee from Fomahand at a time of year when the events actually occur in the following order? In the evening after sunset, as soon as any stars are visible at all, both have already risen, Fohamond admittedly almost half an hour earlier than the Pleiades. But Fomahand rises only a little above the horizon and has already set again long before midnight, while the Pleiades shine high in the sky all night and go below the horizon only at the moment the sun rises. So it seems that if the passage above can be used as a parallel, it would more naturally support Bishop \textsc{Horsley}’s interpretation, which sees in the last two lines the evening setting of the Pleiades coinciding with the morning rising of the Fishes (in spring), were it not for other objections of some weight to that interpretation. [The reviews of the second and third essays and the letters are in Volume IV of these works, among the reviews of writings on subjects in geometry].
